This chapter will focus on presenting the context of this thesis.
In particular in \hyperref[sec:fuzz_background]{ Section 1.1} the foundamentals of fuzzing will be discussed with a particular attention on coverage-guided grey-box fuzzing.
In \hyperref[sec:biniary_rewriting]{Section 1.2} is discussed the background on binary rewriting, its approaches and its difficulties.
\hyperref[sec:problem_statement]{Section 1.3} states the problem addressed by this work, 
and \hyperref[sec:state_of_the_art]{Section 1.4} reviews the state of the art in binary-only fuzzing, focusing also on a description
of AFL/AFL++ internals. In the last section (\hyperref[sec:goals-and-challenges]{Section 1.5}) I will present which are the main 
goals and challenges of this thesis, and why the implemented tool differ from the state of the art .

\section{Fuzzing Background} \label{sec:fuzz_background}
    Fuzzing is nowaday a widely used approach for software testing and vulnerability assesment \cite{Zhu2022,Mallissery2023}. 
    Its main strenght derives from the fact that could rapidly produce various inputs to be fed into a target executable.
    In general a fuzzer can be divided into the \textbf{Input Generator}, that is responsible of provide inputs to the \textbf{Executor} that execute them.
    During the execution the \textbf{Defect Monitor} check if the input executed ends in a crash of the target.\\
    In particular the Input Generator must choose which input provide to the executor between all the inputs in the \textit{Input Space}.
    A subset of the \textit{Input Space} is represented by the \textit{Valid Input Space}.
    This is a subset of the input space, specifically contains all the inputs that are well formed with respect to the 
    application that is being tested.
    For example if the fuzzed target if a pdf-reader, a valid input is represented by a pdf file.
    A good Input Generator must choose inputs that are also valid inputs, because only by doing this there is a probability of find a vulnerability\cite{Zhu2022}.
    So the need of a precise input selection and mutation arise, in order to discover as many edge cases as possible.\\
    In order to do so, the Input Generator must gain some informations from the Executor. 
    Depending on the amount of informations that the Input Generator has, fuzzing can be classified into 
    \textit{White-box}, \textit{Black-box} or \textit{Grey-box}.
    In Black-box fuzzing the fuzzer does not have any knowledge about the execution \cite{Miller1990}.
    In White-box fuzzing instead the source code must always be available and technique of Symbolic and Concolic execution are used \cite{Godefroid2012}.
    In Grey-box fuzzing instead the inputs are choosen using, for example, \textbf{Code Coverage} informations (i.e. number of lines of code visited during execution of the target).
    This approach is called \textit{Coverage-Guided Grey-box fuzzing} and is based on prioritizing the inputs that discover new sections of code.
    Since not always the source code is available the Code Coverage is not always a viable Coverage Measure.
    In general the coverage measures could be of different granularity \cite{CollAFL}. 
    The main three options are:
    \begin{itemize}
        \item \textbf{Block Coverage} : Keep track of the hit count for each \gls{BB}. This approach is used by some 
            popular fuzzers such as libFuzzer or honggfuzz \cite{CollAFL}. This method is often too imprecise because 
            it fails to capture the transitions between bb
        \item \textbf{Edge Coverage} : Keep track of the hit count of every edge in the \gls{CFG} (i.e. transition between two \gls{BB}).
            This is the approach used by AFL and AFL++ \cite{aflplusplus}. Maintain the overheads similar to Block Coverage but could 
            track the transition between \gls{BB} and not only their execution adding some context information.
            Beside this it lack of precision in keep track of the order in which the edges are taken or not.
        \item \textbf{Path Coverage} : This is the most complete metric and keep track of the complete path taken during the execution 
            of a specific input, so keeping track also of the order in which the edges are taken.
            The main drawback of this metric is represented by the fact that a single execution path could become very large, and so 
            incur in large runtime and memory overheads.
    \end{itemize}

    The best approach to Code Coverage, in terms of overhead and effectiveness, is \textbf{Edge Coverage} \cite{CollAFL} which keeps track of which edges 
    in the \gls{CFG} are taken during execution but without the overheads given by keeping track of all the execution path.
    The coverage informations are obtained through the instrumentation of the target application. 
    Depending on the availability of the source code, the instrumentation could be added in different ways:
    \begin{itemize}
        \item The source code is available : In this case the instrumentation is added at \textit{Compile Time}
        \item The source code is not available : In this case there are two different techniques to instrument the application:
        \begin{itemize}
            \item \textbf{\gls{SBI}} : The instrumentation is added before the target is executed but without having the 
                source code, so relying on Static Binary Rewriting techniques.
            \item \textbf{\gls{DBI}} : The instrumentation is added during the program execution, relying on Dynamic Binary Rewriting techniques.
        \end{itemize}
    \end{itemize}
    Indipendently from the availability of the source code, \textit{Coverage-Guided Grey-box fuzzing} technique could be summarized
    as seen in \cref{fig:coverageguidedfuzzing}.
    \begin{figure}[ht]
        \caption{Generic diagram that describe Coverage-Guided Grey-box fuzzing workflow}
        \centering
        \includegraphics[width=\textwidth]{CoverageGuidedGrayBoxFuzzing.png}
        \label{fig:coverageguidedfuzzing}
    \end{figure}
    \\Specifically for \textbf{Edge Coverage} the \cref{fig:coverageguidedfuzzing} could be described as follows:
    \begin{enumerate}
        \item \textbf{Seed Corpus} : Represent the set of inputs on which are computed the mutations. 
            At the beginning is composed of user defined seeds but during fuzzing this set can grow if a specific mutation discover new edges not yet discovered.
        \item \textbf{Input Generator} : Is composed of two main parts:
            \begin{itemize}
                \item \small \textbf{Seed Scheduling} : responsible of decide which seed in the corpus to mutate
                \item \small \textbf{Mutation Scheduling} : responsible of decide which mutation or mutations apply to the selected seed
            \end{itemize}
        \item \textbf{Executor} : Is the component that actually execute the target binary. 
            If the target binary has already been instrumented with Edge Coverage tracking instruction is executed as is, 
            otherwise, in case of \gls{DBI} a dynamic translator must dynamically instrument the binary.
        \item \textbf{Coverage Collection} : In case of Edge Coverage this part is represented by the data structure used 
            to keep track of which edges between \gls{BB} has been taken and which are not.\\
            These information are used, in some tools such as AFL++ \cite{aflplusplus} to guide Seed Scheduling, 
            a technique that, given different inputs with the same level of coverage (i.e. number of edges visited), 
            prioritize the input that is smaller in size and takes less execution time.
        \item \textbf{Defect Monitor} : this component check for exit status of the target to catch any crash or 
            for an eventual timeout of the target or any other unexpected behavior of this.
            If none of these conditions are met, which means that the target ends normally, the coverage is checked 
            to see if the current input had discovered a new edge. 
            In case the input had discovered a new edge, that input is added to the corpus.
        \end{enumerate}
        
\section{Binary Rewriting Background}\label{sec:biniary_rewriting}
        Binary rewriting consist in the alteration af an already compiled executable without having access to the source code 
        while the modified binary remains executable. 
        The process of binary rewriting can be classified into Static Binary Rewriting , if the binary is stored into the 
        persistent memory while rewrited, and Dynamic Binary Rewriting, if the binary is being executed.
        The process of binary rewriting could be divided into four main steps \cite{binaryrewriting}:
        \begin{enumerate}
            \item \textit{Parsing} : In this phase the main goal is to retrieve informations about the binary. 
                In particular the output of this phase is a raw stream of instruction (in machine code) and data.
                These informations are then used in the Analysis step

            \item \textit{Analysis} : The informations gained from the Parsing phase are analyzed to reconstruct the program structure.
            For example the raw stream of machine code is passed through a disassembler that reconstruct the 
            assembly instructions. A structural recory is performed, that include function detection and \gls{CFG} reconstruction, 
            labels and symbols are extracted from the raw stream of data.

            \item \textit{Transformation} : Once the \gls{CFG} has been reconstructed and functions and data boundaries has been recovered, 
                the binary rewriter can modify the binary at the \gls{IP}. These are user defined locations in the binary where 
                the rewriter must add, modify or delete an instruction. The \gls{IP} could be coarse or fine-grained.
                In the first case, for example, the instrumentation could be made at branch level (at the beginning of each \gls{BB}).
                In the second case instead any instruction could be potentially modified.
            
            \item \textit{Code Generation} : In this last step the changes made in the previous steps must be integrated into the 
                binary preserving its executability. This could be done in different ways. For example the changes could be 
                inserted into a new section of the executable and redirect the execution to the new code at the \gls{IP}.
                Otherwise this step could be done also reassemblying a completely new executable starting from the 
                changes made at \gls{IP}. This second approach is much harder because need to preserve all the 
                references inside the excutable and preserve the executable layout.
        \end{enumerate} 
        \begin{figure}[ht]
            \caption{An example of Static Binary Rewriting Workflow}
            \centering
            \includegraphics[width=200pt]{BinaryRewritingScheme.png}
            \label{fig:static_rewrite}
        \end{figure}
        In \cref{fig:static_rewrite} can be seen a general scheme for static binary rewriting.
        In this case all the phases occur prior the execution. So the cost in terms of overheads 
        is paied only one time before the execution. The main problems are disambiguate between code and data 
        (i.e. decide if a stream of bytes represent machine code instructions or are constant placed inline into the .text section) 
        and resolving indirected branches \cite{binaryrewriting,andriesse2016}.
        This problem is in general undecidable and can be resolved only using heuristics. 
        There are three main families of static transformation:
        \begin{itemize}
            \item \textbf{Direct} : Consist in directly modifying the instruction. Works well with RISC systems with instructions having 
                the same lenght. Generally fails to add and remove instruction because this doesn't preserve the offset into the .text section.
            \item \textbf{Minimal-invasive}: The new code is inserted into a new section of the binary and in the \gls{IP} 
                is placed an unconditional jump to the new section. 
                The main liminitation of this approach arise when the new section is far from the \gls{IP} and a jump near (necessary to reach the new section) 
                uses a number of bytes different from the number of bytes used by the instruction at the instrumentation point. 
                Infact in assembly x86-64 a jump short uses 2 bytes and a jump near uses up to 5 bytes.
                So in case at the instrumentation point is needed a jump short for alignment reason, 
                but the new section is too far away and a jump near is needed, it is necessary to add a \textbf{trampoline}.
                This approach can be seen in \cref{fig:trampoline_instrumentation}.
                The tracmpoline based instrumentation is one of the most common approah to binary rewriting \cite{binaryrewriting}.
                
            \item \textbf{Full-translation} : The binary is lifted into an \gls{IR} that makes possible to add an \gls{IP} at 
                instruction level. The \gls{IR} is then reassembled into machine code after the changes.
        \end{itemize}
        \begin{figure}[ht]
            \caption{An example of Static Instrumentation with trampoline technique}
            \centering
            \includegraphics[width=400pt]{trampolineInstrumentation.png}
            \label{fig:trampoline_instrumentation}
        \end{figure}
        With regard to Dynamic Binary Rewriting the Analysis and Transformation phases are executed at runtime for each instruction.
        In this case there are not all the problems that arise in Static Binary Rewriting, since the Binary is being executed, 
        and so the indirected branches could be resolved with the values inside the registers and there is no need for 
        disambiguate between wich are constants inside the .text section and which are instructions.
        The main drawback of this approach consist in prohibitive overheads caused by the dynamic analysis and transformation of each instruction, 
        alongside the fact that only the taken branches could be modified.
        Infact if a branch is not actually taken during execution, its instructions never reach the rewriter.\\
        \\
        A specific type of binary rewriting is the Binary Instrumentation.
        The binary instrumentation uses binary rewriting techniques to add instructions that tracks binary behaviour, 
        log event or, in the case of this thesis, insert edge coverage tracing instructions.
\section{Problem Statement}\label{sec:problem_statement}
    Coverage-Guided Grey-box fuzzing is one of the most effective techniques in automatic vulnerability disclosure on software 
    \cite{zafl,Boehme2016}, 
    due to its capability of generating wide amount of inputs based on the feedback obtained through Code Coverage informations 
    that guide the fuzzer towards unexplored regions of code and therefore find hidden vulnerabilities.
    Infact \textit{Black-box} fuzzing lack of a precise mutation strategy based on execution monitoring while 
    \textit{White-box} fuzzing incur in high overheads given by Symbolic or Concolic execution.
    Furthermore, the latter approach also requires the availability of the source code, not always available.
    So \textit{Grey-box} fuzzing seems to be the most promising trade-off between precision of input mutation and overheads.
    However, also \textit{Grey-box} fuzzing has its own drawback.
    As matter of facts the effectiveness of this approach strongly depends on the availability of the source code.
    Infact both \gls{SBI} and \gls{DBI} have many limitations.
    The main limitation of \gls{DBI} is represented by huge amount of overheads (higher than 600\% overheads 
    in case of AFL++ in QEMU-Mode \cite{zafl}) that makes the exploration of a substantial part 
    of the valid input space almost impossible.
    Instead \gls{SBI} main difficulty arise in modifying already compiled binaries.
    There are several problems regarding Binary Rewriting such as Indirect Control Transfers (e.g. \texttt{call rax}), 
    shift in binary sections or instructions that cause the need of rewrite the instructions that reference these sections and many others 
    as seen in the previous \hyperref[sec:biniary_rewriting]{Section 1.2}.\\
    For all these reasons the preferable technique nowaday is \textit{Compile time Edge Coverage instrumentation} with fuzzers such as 
    AFL++, LibAFL \cite{libafl} o libFuzzer \cite{libfuzzer}.
    However in practice, for industrial secrecy (i.e. proprietary software) or due to precompiled libraries, the applications are 
    often closed source, and so Compile Time Edge Coverage instrumentation become unusable.\\
    In these scenario binary only fuzzing is the only viable approach.
    Although \gls{SBI} presents many challenges, it nonetheless proves to be the most fruitful area, 
    given the low overheads \cite{zafl}.
    In this thesis I will address the problem 
    of statically instrument a binary with Edge Coverage instructions in a way that the instrumented binary is compatible with AFL++,  
    which  has been choosen as fuzzer due to its popularity and because represent the current state of the art implementing 
    most of the recent discovered techniques disclosed demostrated by Andrea Fioraldi et al. \cite{aflplusplus}.\\
    The instrumentation must add as little overheads as possible with respect to compile time overheads and must 
    preserve the correctness of the binary. This is important in fuzzing because less overheads means more executions per second, 
    so a bigger portion of the valid input space could be visited.
   
    \section{State of the Art}\label{sec:state_of_the_art}
        In this section I will disclose some of the current state of the art solutions in terms of binary only greybox fuzzing, 
        differentiating between \gls{SBI} and \gls{DBI} approaches.
        
        \begin{itemize}
            \item \underline{\textbf{\gls{DBI}}}
            \begin{itemize}
                \item \textbf{AFL++ in QEMU/FRIDA mode}\cite{aflplusplus} : AFL++ in QEMU mode uses QEMU in \textit{User-mode Emulation}
                    which consist in translating guest machine code (i.e. guest \gls{BB}) into an intermediate representation 
                    called \gls{TB} and subsequently append the coverage tracing instructions at the end of it.
                    This approach is based on emulating the entire CPU thanks to QEMU virtualized environment. As stated before, 
                    AFL in QEMU mode has overheads in the order of 600\% with respect to the compile time instrumentation \cite{zafl}\\
                    Instead AFL++ in FRIDA mode \cite{aflplusplus-frida-mode} is based on a shared library injection in which is present a frida \cite{frida} component 
                    called \textbf{Stalker} which intercept any \gls{BB} that is about to be executed , copy it into a \textbf{Code Cache} 
                    and dynamically instrument the copy with coverage tracing instructions. 
                    The execution control is then transferred to the copy, execute it, and then re-intercepted by \textbf{Stalker} at the end.
                    Differenctly from QEMU mode doesn't emulate the CPU, but is executed within the memory space of the target program.
                    AFL++ reports that FRIDA mode is approximately 2 to 5 times slower than compile-time instrumentation \cite{aflplusplus_binary_only_targets}.
                
                    \item \textbf{DynamoRIO/Pin}\cite{DynamoRIOPaper,PinPaper} : Both these approaches share the same base principle, called 
                    \textbf{\gls{JIT} code caching}. This consist on take control of the thread when the target application starts, 
                    and from that moment on every \gls{BB} is intercepted befor it is executed and is copied into a private \textbf{Code Cache}.
                    The instrumentation is then added to the copy and the control is transferred to the code cache.
                    The instrumentation process is more or less identical to FRIDA's Stalker mechanism, but are not optimized like taht.
                    Also in this case the implemented version uses AFL\cite{zalewski2017afl} as fuzzer \cite{heuse2018afldynamorio,heuse2018aflpin}.
            \end{itemize}
                \item \underline{\textbf{\gls{SBI}}}
                \begin{itemize}
                    \item \textbf{Retrowrite} \cite{retrowrite}: This tool performs symbolic lifting on the code, replacing every reference
                        discovered during disassembly with a symbolic label. Doing this allows instructions to be added without
                        modifying the references, achieving reassembleable assembly.
                        While RetroWrite can add ASan instrumentation, it cannot natively add AFL-style edge coverage instructions.
                        To do so, the lifted assembly must be instrumented with the afl-gcc tool in order to produce the patched binary.
                        Moreover, community-reported issues \cite{retrowrite_issue28} document missing instrumentation points when
                        RetroWrite's symbolized output is compiled through afl-gcc, caused by a mismatch between RetroWrite's
                        label-naming convention and afl-as's basic-block detection heuristics.
                        Finally, RetroWrite only works with PIE, non-stripped binaries.
                    \item \textbf{AFL-Dyninst} \cite{afldyninst} : AFL-Dyninst employs Dyninst's binary rewriting capabilities 
                            to instrument an existing executable by inserting callbacks at selected basic blocks. 
                            Dyninst's instrumentation mechanism is based on code patching and trampoline redirection \cite{Dyninst}, 
                            which cause many context switches and state saving/restoring.
                            Its performance are slightly better than \gls{DBI} but this approach reaches 500\% of fuzzing overheads 
                            as demostrated by Stefan Nagy et al. \cite{zafl}.
                 
                    \item \textbf{ZAFL} \cite{zafl} : This tool uses two main phases: the first phase consist in 
                        reconstructing an \gls{IR} using the Zipr bynary rewriter. 
                        The obtained \gls{IR} is then used by ZAX to perform an optimization on the \gls{IR}, find the instrumentation points 
                        and in the end instrument these points. ZAFL works inserting calls to a function that update the shared map 
                        at the beginning of each \gls{BB}.
                        Beside many improvements of ZAFL, Jifan Xiao et al. \cite{SystemCallPatternTracing} demostrated that 
                        still cause crashes in binaries in 28\% of the binaries used in the cited paper.
                \end{itemize}
        \end{itemize}
        Before proceding with the next section in which I will talk about the main goals of this thesis, I want to discuss the core functioning of AFL and AFL++, 
        given its importance for understanding my approach.
        \subsection{AFL Functioning} 
        \cite{afl_technical_details}
            As previously stated, AFL and, in particular, AFL++, are among the most widely used grey-box 
            fuzzing frameworks \cite{aflplusplus}.
            AFL++ represent the community based continuation of the initial AFL project.
            The core functioning is virtually the same and consist in a \gls{SHM} region, shared between the fuzzer and the target.
            The \gls{SHM} dimension could vary but the default one is of 65536 bytes.\\
            Each ``cell'' of the  \gls{SHM} contains the number of times a given Edge in the \gls{CFG} has been 
                taken during the execution of the target. Since the \gls{SHM} has 65536 ``cells'', each ``cell'' 
                can contain a number from 0 to 255.
            The ``cell'' position is computed with a specific hash function applied to the IDs of the two \gls{BB} 
            that formed the edge. 
            The main Edge Tracing update rule is summarized in \cref{alg:afl_style_update}. 
            \begin{algorithm}
                \caption{AFL Style SHM update rule}
                \label{alg:afl_style_update}
                \begin{algorithmic}[1]
                    \STATE $idx \gets current\_bb\_id \oplus previous\_bb\_id$ \label{line:xor}
                    \STATE $shared\_map[idx] \gets shared\_map[idx]+1$ \label{line:increment}
                    \STATE $previous\_bb\_id \gets current\_bb\_id \gg 1$ \label{line:shift}
                \end{algorithmic}
            \end{algorithm}
            The line \cref{line:shift} of \cref{alg:afl_style_update} is used to avoid update the same \gls{SHM} 
            location if an edge is traversed forward or backward.
            An example of this could be a simple \texttt{do-while} loop as seen in \cref{alg:do-while} and in \cref{fig:Do-While-CFG}.\\
            \begin{figure}[ht]
                \centering
                \begin{minipage}[h]{0.48\textwidth}
                    \centering
                    \begin{algorithm}[H]
                        \caption{Do-While Example}
                        \label{alg:do-while}
                        \begin{algorithmic}
                            \REPEAT
                                \STATE \COMMENT{Basic Block A }
                            \UNTIL{$Permanence\_Condition$}
                            \STATE \COMMENT{Basic Block C}
                        \end{algorithmic}
                    \end{algorithm}
                Without \cref{line:shift} of \cref{alg:afl_style_update} the hash of Edge1 
                will result in $ID\_A \oplus ID\_B$ and the hash of Edge2 (in case the permanence condition is not met) will result 
                in $ID\_B \oplus ID\_A = ID\_A \oplus ID\_B$ for xor simmetry. But since the hash is the same there will be an hash collision
                that cause the same ``cell'' to be updated two times instead on one.
                \end{minipage}
                \hfill
                \begin{minipage}[h]{0.48\textwidth}
                    \centering
                    \includegraphics[width=100pt]{AFLUpdateAlg3InstructionWhy.png}
                    \caption{Do-while CFG}
                    \label{fig:Do-While-CFG}
                \end{minipage}
            \end{figure}
            The main idea behind AFL is instrumenting each basic block of a target application with the shared map update rule.
            The fuzzer initialize the \gls{SHM} memory with zeros, run the instrumented target that will update the \gls{SHM} 
            with the \cref{alg:afl_style_update},  until it reach the end and then return to the fuzzer.
            The \gls{SHM} ID is comunicated to the target with an environmental variable called \textit{\_\_AFL\_SHM\_ID}.\\
            At this point the fuzzer check the just obtained \gls{SHM}, that il called \texttt{trace\_bits}, searching for any new discovered edge.
            This is done keeping a structure called \texttt{virgin\_bits} that contains all the edges not yet discovered and 
            it is initialized with ones. The compare is executed through $\texttt{trace\_bits} \land \texttt{virgin\_bits}$. If one of the position gives one after the logic AND, 
            this means a new edge has been discovered. In this case the \texttt{virgin\_bits} data structure is updated and  
            the just executed input is saved in the corpus. Otherwise a new input is choosen, mutated and then executed. \\
            This approach requires that for each input a \texttt{fork()} call and then an \texttt{execve()} call are performed, 
            causing a lot of overheads for each input execution.
            Two main optimizations are been implemented:
            \begin{itemize}
                \item \textbf{forkserver}
                \item \textbf{persisten-mode} : \small (vailable only in AFL++)
            \end{itemize}
            \vspace{0.5cm}
            \begin{center}
                \textbf{Forkserver}
            \end{center}
            The forkserver makes it possible to avoid calling \textbf{execve} for every input, performing it only once at the start of execution.
            The forkserver is instrumented before the main of the target application.
            The AFL++ forkserver protocol could be summarized as follows:
            \begin{figure}[ht]
                \caption{AFL++ Forkserver Protocol}
                \centering
                \includegraphics[width=350pt]{AFLPlusPlusForkserverFunctioning.png}
                \label{fig:aflplusplusforkserver}
            \end{figure}
            \footnote{ \small \textit{The differences between AFL forkserver protocol and AFL++ forkserver protocol are mainly the messages exchanged between forkserver and fuzzer}}
            The communication between Fuzzer and Target seen in \cref{fig:aflplusplusforkserver} happens 
            unidirectionally through two different file descriptors created by the fuzzer:
            \begin{itemize}
                \item \textit{FORKSRV\_FD (fd 198)}: used by the forkserver to read a message sent by the fuzzer.
                \item \textit{FORKSRV\_FD + 1 (fd 199)}: used by the forkserver to write a message to the fuzzer.
            \end{itemize}
            \begin{enumerate}
                \item The forkserver sends the ``Welcome'' message to the fuzzer (0x41464c01).
                \item The fuzzer reply with the status code. At this point the forkserver must check if $status \land 0xffffffff = 0x41464c01$.
                        If it is equal continue the execution, otherwise the forkserver stops here.
                \item The forkserver sends the execution options to the fuzzer. 
                    These could indicate the fact that the target is running in persisten mode or other optimizations.
                \item The forkserver sends 0x41464c01 to the fuzzer to end the handshake
                \item The fuzzer write a message to the forkserver (its content is discarded) the handshake is finished.
                            This message is the ``go'' signal for the forkserver.
                \item The forkserver invoke \texttt{fork()}. The parent write the PID of the child to the fuzzer and then invoke \texttt{waitpid()}.\\
                        The child continue the execution of the main until it ends
                \item When the parent exit from the \texttt{waitpid()} function (i.e. the child terminate),  the exit status 
                of the child is sent to the fuzzer and the forkserver waits for another ``go'' from the fuzzer.
            \end{enumerate}        
    \section{Goals and Challenges}\label{sec:goals-and-challenges}
    The main goal of this thesis is to develop a tool capable of inline, static edge-coverage instrumentation
    that preserves binary correctness while achieving performance comparable to compile-time instrumentation.\\
    I have choose to implement the \gls{SBI} approach instead of \gls{DBI} due to its low overheads.
    For doing this I will use GRAPE binary rewrite framework, a novel approach for binary rewriting that preserve binary executability, 
    while introducing small overheads with respect to a trampoline based approach.
    This approach involves several challenges, mainly related to preserving the binary's correctness and behavior.
    Moreover the tool must be compatible with the already existing AFL++ fuzzer. 
    This means implement the forkserver communication protocol inside the target binary without recompiling it.
    This had created not few problems, mainly regarding the preservation of ELF layout after adding new sections.
    With respect to ZAFL approach, the \gls{SHM} update instructions are placed directly inside the \gls{BB}, and not 
    in a function called at each \gls{BB}.
    This approach increment drastically the dimension of the file, causing many instruction repetitions (infact the \gls{SHM} 
    instructions are always the same), but has been chosen because does not cause a context switch at each \gls{BB}.
    Infact, in order to call a function, the state of the registers must be saved and restored. Moreover stack must be allocated, 
    the callee RIP must be saved on the stack alongside the old Base Pointer of the stack.
    All these operations to be performed at each \gls{BB} introduce significant overheads.
    